\section{Arbitrary growth under fixed counting bounds}
\label{sec:arbitrary-growth}

Finiteness of the negative envelope in each compact window does not
by itself impose an exponential growth rate. This remains the case
under both reflection and conjugation symmetry, with fixed linear
origin counts and a fixed logarithmic unit-height bound. The finite
screen calculation supplies the separation between two support
radii needed to construct such an example.

More precisely, let \(G:[2,\infty)\to(0,\infty)\) be any
finite-valued nondecreasing function. There exists a locally finite
divisor \(Z\), with positive integer multiplicities, such that
\begin{gather}
 Z=-\overline Z=\overline Z=-Z,\qquad
 |\Re z|\le1\quad(z\in Z),\qquad
 [1]+[-1]\subset Z,\qquad s_Z=1,
 \label{arb:symmetry}\\
 N_0(Y)\le Y+2\quad(Y\ge0),\qquad
 N_t(1)\le2+\log(2+|t|)\quad(t\in\mathbb R),
 \label{arb:counts}\\
 \cM_Z(L)<\infty\quad(0<L<\infty),\qquad
 \cM_Z(L)\ge G(L)\quad(L\ge2).
 \label{arb:conclusion}
\end{gather}
The symmetries and counts include multiplicity. Here \(s_Z\) is
the supremum of the positive real parts, and
\[
 \cM_Z(L)=
 \sup_{0\ne f\in C_c^\infty((-L,L);\mathbb C)}
 \max\left\{0,-\frac{Q_Z(f,f)}{\|f\|_2^2}\right\}.
\]
The series defining \(Q_Z\) is absolutely convergent on this
test class. For each \(L\ge2\), a negative witness can be
chosen real-valued and either even or odd; no fixed choice of
parity is asserted for all windows. The construction also gives
\begin{equation}\label{arb:weighted-failure}
 \sup_{y\in\mathbb R}
 \sum_{\substack{z=\sigma+i\gamma\in Z\\
                  \sigma>0,\ |\gamma-y|\le1}}
       \nu(z)\sigma^2=\infty.
\end{equation}
No continuity, differentiability, or computability assumption on
\(G\) is used. The conclusion is domination of \(G\), not an
asymptotic formula for \(\cM_Z\).

\subsection{Separating two radii with a finite screen}

We use only the even values \(n\ge2\) of the finite screen
\eqref{scr:screen-definition}, namely
\[
 \mathcal S_{n,\beta}
 =128\sum_{\ell=1}^{n}
   \left[i\beta n\cos\frac{(2\ell-1)\pi}{2n}\right]
       +[\beta]+[-\beta],\qquad \beta>0.
\]
It has total multiplicity \(128n+2\) and is invariant under
both conjugation and reflection. Write
\(e_{n,\beta}(L)=\cM_{\mathcal S_{n,\beta}}(L)\).
The exact finite-screen calculation in
\eqref{scr:envelope} gives, for fixed even \(n\),
\begin{equation}\label{arb:screen-input}
 e_{n,\beta}(L)
   =c_n\beta^{2d}L^{2d+1}
       \bigl(1+O_n((\beta L)^2)\bigr),\qquad
 d=n+1,\qquad c_n>0.
\end{equation}
Its error is not required to be uniform in \(n\).
The finite form is bounded on \(L^2(-L,L)\), with exactly
one negative eigenvalue. Consequently \(e_{n,\beta}(L)>0\),
it is nondecreasing in \(L\), and its value is approached by
unit compact smooth tests. In particular, for each fixed
\(j\ge1\) and even \(n\),
\begin{equation}\label{arb:radius-ratio}
 \frac{e_{n,\beta}(j)}{e_{n,\beta}(2j)}
       \longrightarrow2^{-(2d+1)},\qquad
 e_{n,\beta}(j)\longrightarrow0
       \quad(\beta\downarrow0).
\end{equation}

Put \(A_j=\max\{1,G(2j+2)\}\) and
\(\delta_j=2^{-j}\). Choose an even integer \(n_j\ge2\),
with \(d_j=n_j+1\), sufficiently large that
\begin{equation}\label{arb:degree-choice}
 2^{-(2d_j+1)}<\frac{\delta_j}{32A_j}.
\end{equation}
This requires only that \(A_j\) be finite. Fixing \(n_j\)
before decreasing \(\beta_j\), choose \(\beta_j>0\) so that,
with \(e_j=e_{n_j,\beta_j}\),
\begin{equation}\label{arb:displacement-choice}
 \beta_j\le\min\{2^{-j},1/n_j\},\qquad
 \frac{e_j(j)}{e_j(2j)}\le\frac{\delta_j}{16A_j},\qquad
 e_j(j)\le\frac{\delta_j}{4}.
\end{equation}
Define the integer
\begin{equation}\label{arb:multiplicity}
 m_j=\left\lceil\frac{4A_j}{e_j(2j)}\right\rceil.
\end{equation}
The ceiling is controlled explicitly:
\begin{equation}\label{arb:two-budgets}
 \begin{split}
 m_je_j(j)
 &\le4A_j\frac{e_j(j)}{e_j(2j)}+e_j(j)
   \le\frac{\delta_j}{2},\\
 m_je_j(2j)&\ge4A_j.
 \end{split}
\end{equation}

The displacement can also be chosen to ensure
\begin{equation}\label{arb:weighted-choice}
 m_j\beta_j^2\ge j.
\end{equation}
Indeed, when \(n_j\) is fixed, \eqref{arb:screen-input} gives
\[
 \frac{4A_j\beta^2}{e_{n_j,\beta}(2j)}
 \sim\frac{4A_j}{c_{n_j}(2j)^{2d_j+1}}
                    \beta^{2-2d_j}\longrightarrow\infty.
\]
Since \(d_j\ge3\), this is compatible with all the smallness
requirements in \eqref{arb:displacement-choice}. The order here
is important: impose all sufficiently-small-\(\beta_j\)
requirements first, then define \(m_j\) by
\eqref{arb:multiplicity}; \(m_j\) is not frozen while
\(\beta_j\) is decreased.

Set
\begin{equation}\label{arb:finite-component}
 D_j=m_j\mathcal S_{n_j,\beta_j},\qquad
 W_j=m_j(128n_j+2).
\end{equation}
The real displacements and ordinate offsets of \(D_j\) are
at most one in absolute value. Its total mass \(W_j\) is finite.
By the smooth-test approximation for the finite screen, select
\begin{equation}\label{arb:unmodulated-witness}
 \phi_j\in C_c^\infty((-2j,2j);\mathbb C),\qquad
 \|\phi_j\|_2=1,\qquad Q_{D_j}(\phi_j,\phi_j)<-2A_j.
\end{equation}
For each stage the screen, its multiplicity, and this smooth
test are now fixed. No uniform derivative estimate for the
sequence of tests is needed.

\subsection{A symmetric induction controlling both directions of interference}

For a finite reflected divisor \(D\), define
\begin{equation}\label{arb:absolute-terms}
 T_D(f)=\sum_{z\in D}\nu_D(z)
                  |H_f(z)|\,|H_f(z^\#)|,
 \qquad z^\#=-\overline z.
\end{equation}
Thus \(|Q_D(f,f)|\le T_D(f)\); this quantity controls
individual absolute terms and does not rely on cancellation
within a component. For a fixed compact smooth function \(h\),
integration by parts yields, for every integer \(N\ge1\),
\begin{equation}\label{arb:fourier-decay}
 |H_h(\sigma+it)|\le C_{h,N}(1+|t|)^{-N}
       \quad(|\sigma|\le1,\ t\in\mathbb R).
\end{equation}
For example, apply integration by parts to \(e^{\sigma u}h(u)\);
its derivatives have uniformly bounded integrals for
\(\sigma\in[-1,1]\), and the bounded \(t\)-range is controlled
directly. It follows that, for fixed finite \(D\) in this strip,
\begin{equation}\label{arb:decay-limits}
 T_{D+i\Gamma}(h)\longrightarrow0\quad(|\Gamma|\to\infty),
 \qquad
 T_D(e^{-i\Gamma u}h)\longrightarrow0\quad(|\Gamma|\to\infty).
\end{equation}
These are finite sums with all multiplicities fixed before
\(\Gamma\) tends to infinity.

Take \(D_0=[1]+[-1]\). We recursively choose positive centers
\(\gamma_j\) and put
\begin{equation}\label{arb:symmetric-components}
 C_j^+=D_j+i\gamma_j,\qquad
 C_j^-=D_j-i\gamma_j,\qquad
 E_j=C_j^+\sqcup C_j^-,\qquad
 f_j(u)=e^{-i\gamma_j u}\phi_j(u).
\end{equation}
The notation \(\sqcup\) adds divisor multiplicities. Because
\(D_j=\overline{D_j}\), conjugation exchanges the two
components in \(E_j\). At stage \(j\), require
\begin{equation}\label{arb:center-budgets}
 \begin{gathered}
 \gamma_1\ge10,\qquad
 \gamma_j\ge\gamma_{j-1}+10\quad(j\ge2),\\
 \gamma_j\ge1+2\sum_{i=1}^jW_i,\qquad
 \log\gamma_j\ge W_j,
 \end{gathered}
\end{equation}
as well as
\begin{equation}\label{arb:new-on-old}
 T_{E_j}(f_k)\le2^{-j-k-2}\quad(1\le k<j),
\end{equation}
and
\begin{equation}\label{arb:old-on-new}
 T_{D_0\sqcup E_1\sqcup\cdots\sqcup E_{j-1}}(f_j)
                   +T_{C_j^-}(f_j)\le\frac14.
\end{equation}

All these requirements can be satisfied simultaneously.
The counting conditions in \eqref{arb:center-budgets} impose
only a finite lower threshold on \(\gamma_j\). Each term in
\eqref{arb:new-on-old} tends to zero by the first limit in
\eqref{arb:decay-limits}, applied to the two translated copies
of the fixed \(D_j\) and each fixed old test. The first term
of \eqref{arb:old-on-new} tends to zero by the second limit,
with the old finite divisor fixed. The current opposite-height
copy obeys the exact identity
\begin{equation}\label{arb:opposite-height}
 T_{C_j^-}(f_j)
     =T_{D_j}(e^{-2i\gamma_j u}\phi_j)
                   \longrightarrow0\quad(\gamma_j\to\infty).
\end{equation}
Thus this last interaction has frequency separation
\(2\gamma_j\). There are only finitely many constraints at
each stage, and \(f_j\) depends on the new center only through
modulation of the already fixed \(\phi_j\). This proves the
induction without a simultaneous infinite choice or a uniform
bound on the smoothing constants.

Define
\begin{equation}\label{arb:divisor}
 Z=D_0\sqcup\bigsqcup_{j\ge1}E_j.
\end{equation}
Every component is finite, its centers tend to infinity, and
the enclosing bands \([\pm\gamma_j-1,\pm\gamma_j+1]\)
are disjoint. Therefore \(Z\) is locally finite. All
multiplicities are positive integers. Imaginary translation
preserves reflection since
\((z+i\gamma)^\#=z^\#+i\gamma\); conjugation exchanges
the two bands of each \(E_j\). This proves all the
symmetries in \eqref{arb:symmetry}. The strip condition follows
from the displacement choices, and the base pair realizes
\(s_Z=1\).

\subsection{Counts, including partially entered bands}

Fix \(Y\ge0\). If there is an index with
\(\gamma_j-1\le Y\), let \(k\) be the largest such index.
Every component contributing to \(N_0(Y)\) has index at most
\(k\). Each of its two bands contributes at most its full
mass \(W_j\), even if the interval has entered only part of
the band. Consequently
\begin{equation}\label{arb:origin-count}
 N_0(Y)\le2+2\sum_{j=1}^kW_j
       \le2+\gamma_k-1\le Y+2.
\end{equation}
If no such index exists, only the base pair can contribute,
and the same bound holds. The use of the first edge
\(\gamma_k-1\), rather than the center of a fully entered
band, proves this at all endpoints as well.

A closed interval \([t-1,t+1]\) can meet at most one of the
signed enclosing bands, since their centers are separated by
at least ten. It cannot meet both such a band and the base
ordinate zero. If it meets a band of index \(j\), then
\(|t|\ge\gamma_j-2\), so
\begin{equation}\label{arb:local-count}
 N_t(1)\le W_j\le\log\gamma_j\le\log(2+|t|).
\end{equation}
If it meets only the base ordinate, its count is at most two;
otherwise its count is zero. This proves \eqref{arb:counts}
with constants independent of \(G\). The factor two in the
origin-count budget accounts for both conjugate bands; it is
not needed in the local count, whose interval meets at most
one band.

\subsection{Absolute convergence and local semiboundedness}

The origin count implies
\begin{equation}\label{arb:inverse-ordinate-sum}
 \sum_{z\in Z}\nu(z)(1+|\Im z|)^{-2}<\infty.
\end{equation}
Indeed the bounded central region is finite, and the part
with \(2^q<|\Im z|\le2^{q+1}\) is bounded by a constant
times \(2^{-q}\). For any fixed compact smooth \(f\),
\eqref{arb:fourier-decay} bounds the absolute divisor summand
by \(C_f\nu(z)(1+|\Im z|)^{-2}\). Hence the series for
\(Q_Z(f,f)\) converges absolutely, and may be grouped by
\eqref{arb:divisor}. The same reasoning applies to the
sesquilinear form on two compact smooth tests.

Fix a finite radius \(L>0\), and set
\(K=\max\{1,\lceil L\rceil\}\). For \(j\ge K\), each
of the two imaginary translations in \(E_j\) has negative
envelope \(m_je_j(L)\). This follows from
\[
 H_{e^{i\gamma u}f}(z)=H_f(z+i\gamma),
\]
because modulation preserves the norm and support. Since
\(L\le j\), monotonicity and \eqref{arb:two-budgets} imply
\begin{equation}\label{arb:tail-lower}
 Q_{E_j}(f,f)\ge-2m_je_j(L)\|f\|_2^2
              \ge-\delta_j\|f\|_2^2
 \quad\bigl(f\in C_c^\infty((-L,L))\bigr).
\end{equation}
Thus the doubled negative budget of all sufficiently late
components is summable.

The finitely many earlier components have finite negative
envelopes. For the base pair, the even and odd features give
\[
 Q_{D_0}(f,f)=
 2|\langle f,\cosh u\rangle|^2
 -2|\langle f,\sinh u\rangle|^2,
\]
whose negative envelope is
\[
 e_0(L)=2\|\sinh u\|_{L^2(-L,L)}^2
                         =\sinh(2L)-2L.
\]
Summing the lower bounds for finite unions and then using the
absolute convergence just established yields
\begin{equation}\label{arb:finite-bottom}
 Q_Z(f,f)\ge-B_L\|f\|_2^2,
 \qquad
 B_L=e_0(L)+2\sum_{j<K}m_je_j(L)+\sum_{j\ge K}2^{-j}<\infty.
\end{equation}
This proves the first assertion in \eqref{arb:conclusion}
for every finite \(L\), including \(L<1\). No uniform
convergence over the unit ball is required: the negative
lower bounds are uniform and summable, while the original
series converges absolutely separately on each smooth test.
In particular, we do not assert that the full positive axis
part is a bounded operator on \(L^2(-L,L)\), or that a
closed realization of the full infinite form has been
constructed.

\subsection{Survival of the prescribed witnesses}

On its designated positive-height component the unit smooth
test \(f_j\) satisfies exactly
\begin{equation}\label{arb:own-component}
 Q_{C_j^+}(f_j,f_j)=Q_{D_j}(\phi_j,\phi_j)<-2A_j.
\end{equation}
All old components, the base pair, and the current
negative-height partner together cost at most \(1/4\) in
absolute value by \eqref{arb:old-on-new}. For each future
stage \(k>j\), its constraint \eqref{arb:new-on-old} gives
\(T_{E_k}(f_j)\le2^{-k-j-2}\). Therefore
\begin{equation}\label{arb:future-tail}
 \sum_{k>j}T_{E_k}(f_j)
 \le\sum_{k>j}2^{-k-j-2}=2^{-2j-2}\le\frac1{16}.
\end{equation}
Using absolute convergence and \(A_j\ge1\), we obtain
\begin{equation}\label{arb:surviving-witness}
 Q_Z(f_j,f_j)<-2A_j+\frac14+\frac1{16}<-A_j.
\end{equation}
The test is supported strictly inside \((-2j,2j)\).
For any \(L\ge2\), take \(j=\lfloor L/2\rfloor\ge1\).
Then \(2j\le L<2j+2\), and the increasing test classes and
monotonicity of \(G\) give
\[
 \cM_Z(L)\ge\cM_Z(2j)\ge A_j
                  \ge G(2j+2)\ge G(L).
\]
This proves domination at every radius \(L\ge2\), not merely
along a sparse sequence.

We finish by checking the additional assertions. At height
\(\gamma_j\), the positive-real target
\(\beta_j+i\gamma_j\) has multiplicity \(m_j\), so its
contribution to the sum in \eqref{arb:weighted-failure} is
\(m_j\beta_j^2\ge j\). This proves that assertion.

For real tests \(u,v\), conjugation symmetry of \(Z\) makes
\(Q_Z(u,v)\) real. Together with Hermitian symmetry this gives
\[
 Q_Z(u+iv,u+iv)=Q_Z(u,u)+Q_Z(v,v),\qquad
 \|u+iv\|_2^2=\|u\|_2^2+\|v\|_2^2.
\]
Thus a complex negative test has a nonzero real or imaginary
part with at least its negative Rayleigh margin. Invariance
under \(z\mapsto-z\) makes the form invariant under
\(Rf(u)=f(-u)\); its even and odd real subspaces are
orthogonal for both the form and the norm. Splitting the
chosen real test into those two parts therefore preserves
the margin in at least one nonzero part. Normalization gives
a real, definite-parity unit witness, without changing its
strict compact support.

For instance, the choice \(G(L)=e^{L^2}\) produces a divisor
satisfying \eqref{arb:symmetry}--\eqref{arb:counts}, with a
finite bottom in every finite window, but
\[
 \lim_{L\to\infty}\frac{\log\cM_Z(L)}{L}=+\infty.
\]
Hence local semiboundedness and these fixed count bounds do
not imply a finite exponential growth rate, even with full
reflection and conjugation symmetry. The examples use
multiplicities and are abstract divisors. No simple-divisor
assertion, arithmetic realization, Euler product, or
Riemann--von Mangoldt asymptotic is imposed.
